%! Multiplying \& Dividing Rational Expressions — Unit 5, Lesson 5.2 — Guided Notes (Answer Key)
\documentclass[10pt]{article}
\usepackage{algebra2-article}
\usepackage{algebra2-key}   % pulls in -boxes; do NOT also load -boxes
\newcommand{\TallMath}[1]{$\displaystyle #1\rule[-1.4em]{0pt}{3.2em}$}
\newcommand{\lgrid}[4]{%
  \draw[step=1, linegray!40, very thin] (#1,#3) grid (#2,#4);
  \draw[->, semithick, charcoal] (#1,0)--(#2,0) node[below right, font=\scriptsize]{$x$};
  \draw[->, semithick, charcoal] (0,#3)--(0,#4) node[left, font=\scriptsize]{$y$};}
% Vocab answer for the key: bold forest term, then the filled definition on a write-line.
% The leading and trailing \par are required: \ansline ends with \dotfill but does NOT end the
% paragraph, so without them each term label gets pulled onto the previous answer's line.
\newcommand{\vocabans}[2]{%
  \par\noindent\textbf{\textcolor{forest}{#1:}}\\[1pt]\ansline{#2}\par}

\begin{document}

\pageheader{Unit 5, Lesson 5.2}{Guided Notes --- Answer Key}
\namedateperiod

\vspace{0.08in}

\begin{objectivebox}
\textbf{By the end of this lesson, I will be able to\dots}
\begin{itemize}\setlength{\itemsep}{2pt}
  \item multiply rational expressions by \emph{factoring first} and dividing out common factors across the whole product --- never by multiplying out;
  \item divide rational expressions by multiplying by the \emph{reciprocal} of the divisor (keep--change--flip), flipping \emph{before} cancelling anything;
  \item list \emph{every} restriction from all \textbf{three} sources --- each original denominator \emph{and} the numerator of the divisor.
\end{itemize}
\end{objectivebox}

\vspace{0.05in}

\begin{vocabbox}
\small
Fill in each term as we build it together.
\par\vspace{2pt}   % \par is required: \vocabans's \noindent is a no-op mid-paragraph
\vocabans{Reciprocal (multiplicative inverse)}{the fraction turned upside down: $\frac{c}{d}$ and $\frac{d}{c}$, whose product is $1$ --- and it exists only when $c\neq0$}
\vocabans{Divisor}{the expression you are dividing \emph{by} --- the fraction after the $\div$ sign}
\vocabans{Keep--change--flip}{keep the first fraction, change $\div$ to $\times$, flip the divisor: $\frac ab\div\frac cd=\frac ab\cdot\frac dc$}
\vocabans{Hidden restriction}{an excluded value that appears in neither the simplified answer nor (for a divisor's numerator) the original denominators --- but is excluded all the same}
\end{vocabbox}

\begin{hookbox}
Two expressions, one table --- again. Let
\[
  Q(x)=\frac{x+1}{x-5}\div\frac{x-2}{x+3}
  \qquad\text{and}\qquad
  R(x)=\frac{(x+1)(x+3)}{(x-5)(x-2)} .
\]
$R$ is what you get when you keep--change--flip $Q$ and multiply across.
\begin{center}
\renewcommand{\arraystretch}{1.35}
\begin{tabular}{|l|c|c|c|c|}
\hline
$x$ & $0$ & $1$ & $2$ & $-3$ \\ \hline
$Q(x)$ & \ans{$\frac{3}{10}$} & \ans{$2$} & \ans{undef.} & \ans{undef.} \\ \hline
$R(x)$ & \ans{$\frac{3}{10}$} & \ans{$2$} & \ans{undef.} & \ans{$0$} \\ \hline
\end{tabular}
\end{center}
\begin{itemize}\setlength{\itemsep}{1pt}
  \item Both are undefined at $x=2$ --- but for the \emph{same} reason? \ansline{No. In $R$ the factor $(x-2)$ is a denominator; in $Q$ it is on \emph{top} of the divisor, and it makes the whole divisor $0$ --- so $Q$ is dividing by zero.}
  \item There is exactly one column where they disagree. Which one, and what breaks in $Q$ that does
        not break in $R$? \ansline{At $x=-3$: the divisor $\frac{x-2}{x+3}$ is undefined, so $Q(-3)$ does not exist --- but flipping moved $(x+3)$ upstairs, so $R(-3)=0$.}
\end{itemize}
Yesterday the danger was that cancelling hides a restriction. Today there is a second danger:
\textbf{flipping hides a different one.} The safe habit is to read the restrictions off the problem
\emph{as it was written} --- and then add one more.
\end{hookbox}

\begin{notesbox}{1. Multiplying: the Rule Is Easy, the Discipline Is Factoring}
For numbers and for polynomials, the rule is the same --- multiply straight across:
\[
  \frac{a}{b}\cdot\frac{c}{d}=\frac{a\,c}{b\,d}, \qquad b\neq0,\ d\neq0 .
\]

\vspace{3pt}
\textbf{But you should almost never actually multiply across first.} In the Warm-Up,
$\dfrac{6}{35}\cdot\dfrac{25}{9}$ took one line by factoring and dividing out; multiplying first gives
$\dfrac{150}{315}$, which you then have to reduce anyway. With polynomials it is worse than tedious:
an expanded product \emph{hides its own factors}, and Unit~4 showed you exactly how much work getting
them back costs.

\vspace{3pt}
\begin{center}
\fbox{\parbox{0.88\linewidth}{\centering\vspace{2pt}
\textbf{Factor everything first. Divide out. Multiply last --- and often not at all.}\vspace{2pt}}}
\end{center}

\vspace{4pt}
\textbf{Monomial factors first.} Simplify $\dfrac{6x^2}{5y}\cdot\dfrac{15y^3}{8x^5}$.
\begin{itemize}[leftmargin=1.5em, itemsep=3pt]
  \item Restrictions come from \emph{both} denominators: $5y\neq0$ gives $y\neq$ \ans{0}, and
        $8x^5\neq0$ gives $x\neq$ \ans{0}.
  \item Divide out the numbers ($6$ and $8$ share \ans{2}; $5$ and $15$ share \ans{5})
        and the shared powers. Answer: \ans{$\frac{9y^2}{4x^3}$}.
\end{itemize}
\end{notesbox}

\begin{notesbox}{2. The Four Steps for a Product}
\begin{enumerate}[leftmargin=1.9em, itemsep=2pt]
  \item \textbf{Factor} every numerator and every denominator completely.
  \item \textbf{List the restrictions} from \emph{every} original denominator --- there are two of them now.
  \item \textbf{Divide out} any numerator factor against any denominator factor --- across the \emph{whole} product, not just within one fraction.
  \item \textbf{Write} the result as a single fraction, restrictions attached. Leave it \emph{factored}.
\end{enumerate}

\vspace{4pt}
\textbf{The anchor, worked.} \; $\dfrac{x^2-9}{x^2+2x-8}\cdot\dfrac{x+4}{x^2+3x}$
\begin{itemize}[leftmargin=1.5em, itemsep=3pt]
  \item Factor all four parts. \; $x^2-9=$ \ans{$(x-3)(x+3)$}; \quad $x^2+2x-8=$ \ans{$(x+4)(x-2)$};
        \quad $x^2+3x=$ \ans{$x(x+3)$}.
  \item Restrictions from the \emph{first} denominator: $x\neq$ \ans{$-4,\ 2$}; \; from the
        \emph{second}: $x\neq$ \ans{$0,\ -3$}.
  \item Divide out \ans{$(x+4)$} and \ans{$(x+3)$}. \; Result: \ans{$\frac{x-3}{x(x-2)}$}, \;
        $x\neq$ \ans{$0,\ 2,\ -3,\ -4$}.
\end{itemize}

\vspace{3pt}
Notice what carried over from Lesson~5.1 and what is new. Same as yesterday: two of those four
restrictions are now \textbf{invisible} in the answer. New today: a cancelled factor can come from
\emph{either} fraction, so you have two denominators to read before you touch anything.

\vspace{5pt}
\textbf{Opposite binomials still count} (Lesson 5.1). Simplify
$\dfrac{x^2-25}{x+2}\cdot\dfrac{2x+4}{5-x}$.
\begin{itemize}[leftmargin=1.5em, itemsep=3pt]
  \item $5-x=$ \ans{$-1$} $\cdot\,(x-5)$, so the $(x-5)$'s divide out and leave a factor of
        \ans{$-1$}.
  \item Answer: \ans{$-2(x+5)$}, \; $x\neq$ \ans{$-2,\ 5$}.
\end{itemize}
\end{notesbox}

\begin{notesbox}{3. Dividing: Keep, Change, Flip --- and Why}
The \textbf{reciprocal} of $\dfrac{c}{d}$ is $\dfrac{d}{c}$, and the two multiply to
\ans{1}. Dividing by an expression means multiplying by whatever \emph{undoes} it, so:
\[
  \frac{a}{b}\div\frac{c}{d} \;=\; \frac{a}{b}\cdot\frac{d}{c} .
\]
\textbf{Keep} the first fraction, \textbf{change} $\div$ to $\times$, \textbf{flip} the second.

\vspace{4pt}
\textbf{One warning worth a whole page: flip \emph{first}, then cancel.} Watch what happens if you
cancel across the $\div$ sign:
\begin{center}
\renewcommand{\arraystretch}{1.4}
\begin{tabularx}{\linewidth}{@{}X|X@{}}
\hline
\textbf{Cancelled first (wrong)} & \textbf{Flipped first (right)} \\
\hline
$\dfrac{4}{3}\div\dfrac{5}{4}\;\to\;\dfrac{1}{3}\div\dfrac{5}{1}=$ \ans{$\frac{1}{15}$} &
$\dfrac{4}{3}\div\dfrac{5}{4}=\dfrac{4}{3}\cdot\dfrac{4}{5}=$ \ans{$\frac{16}{15}$} \\
\hline
\end{tabularx}
\end{center}
The two answers are not equal, so cancelling across a $\div$ is \ans{illegal}. The $\div$ symbol is
not a fraction bar --- until you flip, nothing is lined up yet.

\vspace{3pt}
\textbf{Also worth naming:} $\dfrac{c}{d}$ has a reciprocal only when $c\neq0$. That one sentence is
the whole new idea of Section~4.
\end{notesbox}

\begin{notesbox}{4. The Three Sources of a Restriction}
Take any quotient $\;\dfrac{A}{B}\div\dfrac{C}{D}$. Three separate things must be true:
\begin{center}
\renewcommand{\arraystretch}{1.5}
\begin{tabularx}{\linewidth}{@{}c l X@{}}
\hline
\textbf{\#} & \textbf{Condition} & \textbf{Why --- and where it hides} \\
\hline
1 & $B\neq0$ & the first fraction has to exist. \emph{Visible} in the problem and (usually) in the answer. \\
2 & $D\neq0$ & the divisor has to exist before you can divide by it. Goes \textbf{invisible after the flip} --- $D$ moves upstairs. \\
3 & $C\neq0$ & you cannot divide by \ans{0}, and the divisor equals zero exactly when its \emph{numerator} does. \textbf{Invisible before the flip} --- $C$ is upstairs in the problem as written. \\
\hline
\end{tabularx}
\end{center}

\vspace{3pt}
\begin{center}
\fbox{\parbox{0.92\linewidth}{\centering\vspace{2pt}
\textbf{Flipping trades one blind spot for another.} Write the restrictions from the problem
\emph{as written}, then add the \textbf{divisor's numerator}.\vspace{2pt}}}
\end{center}

\vspace{4pt}
\textbf{Back to the Hook.} \; $Q(x)=\dfrac{x+1}{x-5}\div\dfrac{x-2}{x+3}$
\begin{itemize}[leftmargin=1.5em, itemsep=3pt]
  \item Source 1 (first denominator): $x\neq$ \ans{5}. \quad Source 2 (divisor's denominator):
        $x\neq$ \ans{$-3$}.
  \item Source 3 (divisor's \emph{numerator}): $x\neq$ \ans{2}.
  \item Simplified: \ans{$\frac{(x+1)(x+3)}{(x-5)(x-2)}$}, \; $x\neq$ \ans{$5,\ 2,\ -3$}.
\end{itemize}
That is why the table disagreed at $x=-3$: there $Q$ is being asked to divide by
$\dfrac{-5}{0}$, which is not a number at all --- while $R$, where the $(x+3)$ has quietly moved to the
top, is perfectly happy.

\vspace{5pt}
\textbf{The division anchor, worked.} \; $\dfrac{x^2-4}{x^2+7x+12}\div\dfrac{x^2+2x}{x+3}$
\begin{itemize}[leftmargin=1.5em, itemsep=3pt]
  \item Factor first. \; $x^2-4=$ \ans{$(x-2)(x+2)$}; \quad $x^2+7x+12=$ \ans{$(x+3)(x+4)$}; \quad
        $x^2+2x=$ \ans{$x(x+2)$}.
  \item Restrictions, all three sources: $x\neq$ \ans{$-3,\ -4$ (source 1), $-3$ (source 2), $0,\ -2$ (source 3)}
  \item Keep--change--flip, divide out, and write the answer: \ans{$\frac{x-2}{x(x+4)}$}
\end{itemize}
Two of those four restrictions do not appear anywhere in your final answer. They still belong to it.
\end{notesbox}

\begin{notesbox}{5. The Picture: a Product With Two Holes}
Every restriction you list is a point the graph does not get to keep. Consider
\[
  P(x)=\frac{x^2-1}{x}\cdot\frac{x}{x+1} .
\]
\begin{itemize}[leftmargin=1.5em, itemsep=3pt]
  \item Factored: $\dfrac{(x-1)(x+1)}{x}\cdot\dfrac{x}{x+1}$. \; Restrictions: $x\neq$ \ans{$0,\ -1$}.
  \item Everything divides out except \ans{$x-1$} --- so away from those two inputs, $P$ \emph{is}
        a straight line.
\end{itemize}
\begin{center}
\begin{tikzpicture}[scale=0.46]
  \lgrid{-4}{5}{-5}{4}
  \begin{scope}
    \clip (-4,-5) rectangle (5,4);
    \draw[very thick, forest] (-4,-5)--(5,4);
  \end{scope}
  \draw[forest, very thick, fill=white] (0,-1) circle (4.2pt);
  \draw[forest, very thick, fill=white] (-1,-2) circle (4.2pt);
  \node[charcoal, font=\scriptsize] at (-2.3,2.6) {two holes};
  \draw[charcoal, thin, ->] (-2.0,2.2) -- (-0.2,-0.75);
  \draw[charcoal, thin, ->] (-2.0,2.2) -- (-1.1,-1.7);
\end{tikzpicture}
\end{center}
\begin{enumerate}[leftmargin=1.7em, itemsep=3pt]
  \item The hole at $x=0$ came from the denominator of the \ans{first} fraction; the hole at
        $x=-1$ came from the denominator of the \ans{second} fraction.
  \item Read the holes' heights off the \emph{simplified} form $y=x-1$: they sit at \ans{$(0,-1)$}
        and \ans{$(-1,-2)$}.
  \item Neither excluded value shows up in the simplified expression $x-1$ --- there is no denominator
        left to show them. Explain why they are still part of the answer. \ansline{because the restrictions belong to the \emph{original} product: at $x=0$ and $x=-1$ one of the two fractions has no value, so $P$ has none either --- the line $y=x-1$ only tells you where the graph was headed}
\end{enumerate}
\end{notesbox}

\begin{practicebox}
Simplify completely and state \emph{all} restrictions. Show the factored line.

\vspace{3pt}
\begin{enumerate}[leftmargin=1.7em, itemsep=8pt]
  \item $\dfrac{x^2-x-6}{x^2-1}\cdot\dfrac{x+1}{x-3}$ \hfill Factored: \ans{$\frac{(x-3)(x+2)}{(x-1)(x+1)}\cdot\frac{x+1}{x-3}$} \\[4pt]
        Simplest form: \ans{$\frac{x+2}{x-1}$} \quad Restrictions: $x\neq$ \ans{$1,\ -1,\ 3$}
  \item $\dfrac{x^2-16}{3x}\div\dfrac{x+4}{6x^2}$ \hfill Factored (after the flip): \ans{$\frac{(x-4)(x+4)}{3x}\cdot\frac{6x^2}{x+4}$} \\[4pt]
        Simplest form: \ans{$2x(x-4)$} \quad Restrictions: $x\neq$ \ans{$0,\ -4$}
  \item Your answer to problem 2 has no denominator at all. Explain why it still carries restrictions,
        and name which source each one came from.
        \ansline{$x\neq0$ comes from the original denominators $3x$ and $6x^2$ (sources 1 and 2); $x\neq-4$ comes from}
        \ansline{the divisor's \emph{numerator} $x+4$ (source 3) --- at $x=-4$ the problem asks you to divide by zero.}
\end{enumerate}
\end{practicebox}

\begin{teachernote}
\textbf{Pacing.} Sections 1--3 are the mechanics and go quickly if yesterday landed; \textbf{Section 4
is the lesson}. If time runs short, move the opposite-binomial example in Section~2 to homework rather
than cutting Section~4 or 5 --- the three-source rule is what makes today more than ``5.1 with two
fractions.''

\vspace{3pt}
\textbf{The one sentence to repeat.} \emph{Flipping trades one blind spot for another.} Students accept
the divisor-numerator restriction in the Warm-Up (item 1c, with numbers) and forget it the moment
letters appear. When they do, send them back to $\frac49\div\frac{n}{15}$, not to a rule.

\vspace{3pt}
\textbf{Errors to expect.} (1) Multiplying out first, producing a quartic nobody wants to re-factor ---
this is a habit, so name it early and often. (2) Cancelling \emph{across} the $\div$ before flipping;
the Section~3 table is the disproof, so make students compute both cells. (3) Flipping the
\emph{first} fraction instead of the divisor. (4) Missing the divisor's-numerator restriction --- the
signature error of the day, and the one the exit ticket's MC item is built to catch. (5) Reading the
restrictions off the final answer, which is yesterday's signature error and still costs a point today.

\vspace{3pt}
\textbf{Section 5} is worth protecting. A product of two rational expressions whose graph is a plain
line with two holes punched in it is the clearest possible picture of why restrictions survive
cancelling --- and it is the same picture 5.5 will build from an equation.
\end{teachernote}

\end{document}
